\documentclass[../main.tex]{subfiles}
\begin{document}
%==========================================TIÊU ĐỀ TẬP========================================
\begin{table}[h]
  \begin{adjustwidth}{-1cm}{}
    \begin{tabular}{>{\centering\arraybackslash}p{6cm}|>{\raggedright\arraybackslash}p{12.5cm}}
      \multirow{5}{6cm}
      % Cột bên trái: ÉPISODE và số tập
      {\\[-40pt]
        \flaregothic\fontsize{20pt}{20pt}\selectfont \centering
        \color{eptype}ÉPISODE\color{black}\\
        \burbank\fontsize{72pt}{72pt}\selectfont
        \color{epnum}69\color{black}\\[6pt]
        \cafeteria\fontsize{20pt}{20pt}\selectfont\color{epnum}(5-13)\color{black}
        \\[18pt]
      }
       &
      % Dòng thứ nhất: tên tập tiếng Nhật
      {\fotskip\fontsize{18pt}{18pt}\selectfont \ruby{全}{ぜん}\ruby{部}{ぶ}\ruby{水}{みず}に\ruby{流}{なが}せ！\ruby{流}{なが}しそうめん\ruby{バ}{ば}\ruby{ト}{と}\ruby{ル}{る}！
      }                            \\[6pt]
      % Dòng thứ hai: tên tập tiếng Anh
       & {\cafeteria\fontsize{18pt}{18pt}\selectfont Nagashi Somen}                                                                                                             \\[4pt]
      % Dòng thứ ba: tên tập tiếng Việt
       & {\vnmsans\fontsize{18pt}{18pt}\selectfont Mì trôi ống trúc}                                                                                                            \\[8pt]
      % Dòng thứ tư: Nhân vật xuất hiện (theo đúng thứ tự debut)
       & {\caviar{\fontsize{14pt}{14pt}\selectfont Nhân vật: \emp{kuon1} \emp{ryuuhou1} \emp{fubuki} \emp{haato} \emp{ayame} \emp{okayu} \emp{korone} \emp{kanata} \emp{coco}}}
      \\[8pt]
      % Dòng cuối cùng: Ngày phát hành
       & {\continuum\fontsize{16pt}{16pt}\selectfont Ngày phát hành: 30/08/2020}                                                                                                \\[18pt]
    \end{tabular}
  \end{adjustwidth}
\end{table}
%=========================================TRUYỆN CHÍNH========================================
\color{black}

% Hộp tóm tắt
\txtbx[2]{{\color{vol} \uline{Tóm tắt:}} Những ngày cuối cùng của mùa hè, Kuon cùng Ayame, Korone, Okayu, Fubuki và Ryuuhou tổ chức tiệc mì trôi ống trúc trên bãi biển. Hai con tôm được thả xuống ống nước và chạy đua ra biển, và Okayu chiến thắng. Khi Fubuki phát hiện quên mang mì, Haato xuất hiện với một xe bán mì somen và đòi giá cắt cổ, nhưng cả nhóm mặc cả thành công. Sau đó, hóa ra toàn bộ là một đoạn phim giáo dục do Coco làm để dạy bài học ``tránh mua hàng qua trung gian'', khiến Kanata hoài nghi về tính thực tế. Kết thúc với hậu trường ghi hình và lời hứa về một bữa tiệc cà ri cuối tháng Tám.}

\pov{Hikawa Kuon}

Ngày áp chót của tháng Tám, cái nóng oi ả vẫn còn vương vấn như một lời nhắc nhở rằng mùa hè sắp khép lại. Những tia nắng cuối cùng vẫn còn rực rỡ, nhưng đã bắt đầu dịu đi, mang theo một chút gì đó bâng khuâng. Trước khi bước sang tháng Chín, chúng tôi quyết định tổ chức một bữa liên hoan để đánh dấu những ngày cuối cùng của mùa hè năm nay, một cái cớ hoàn hảo để tạm gác công việc và tận hưởng không khí thoải mái.

Ngồi trong văn phòng, tôi đang thảo luận với một vài người thì cánh cửa bật mở. Ryuuhou lao vào như một cơn gió, hớn hở sán tới tôi: ``Onii-san! Em nghe nói có tiệc à!? Cho em tham gia với!''

Tôi nhướng mày, giọng đầy nghi ngờ: ``Em nghe ở đâu thế? Mới chỉ có ý tưởng thôi mà?''

Cô em gái cười tinh quái, nháy mắt một cái: ``Em có tai thính mà. Mà dù sao thì em cũng muốn tham gia! Ở nhà chán lắm, trời lại còn nóng nữa.''

Tôi thở dài, nhưng không thể từ chối. Đúng hơn là nếu từ chối thì em ấy sẽ lại viện cớ đủ lý do trên trời dưới biển chỉ ăn ké; cái này tôi biết thừa rồi.

``Mình mở tiệc mì trôi ống trúc đi!'' Ayame hào hứng đề xuất, đôi mắt sáng rực chẳng khác gì những ngôi sao đêm hè. Cô vung tay lên, như thể đang vẽ nên một bức tranh mùa hè hoàn hảo.

Tôi, Korone và Okayu đưa mắt nhìn nhau. Ý tưởng này thực sự không tệ chút nào. Một bữa tiệc ngoài trời, nước mát lạnh, mì trôi theo dòng chảy của ống trúc\dots\ nghe thôi đã thấy thú vị, một trải nghiệm mùa hè đúng chất Nhật Bản.

Nàng Miêu gật đầu ngay lập tức, giọng đều đều nhưng đầy hào hứng: ``Được thôi. Để tôi sẽ mang theo một ít đồ ăn kèm nha\~{}''

Ryuuhou vỗ tay, mắt sáng rực như một đứa trẻ vừa được hứa đi chơi: ``Tuyệt quá! Em chưa từng ăn mì trôi ống trúc bao giờ! Nhìn trong anime thấy ngon lắm!''

Tôi không nói gì khác, ngoài việc chỉ hùa theo, vì biết rằng đã quyết định rồi thì không thể thay đổi: ``Được chứ. Không tồi. Nhưng phải chuẩn bị kỹ càng đấy.''

\img{5-13a}

Nàng Khuyển giơ tay hăng hái, đôi mắt long lanh như một chú chó vừa thấy xương: ``Vậy tôi lo phần ăn uống nhé! Tôi sẽ mang nhiều đồ ăn ngon.''

Okayu cũng hùa theo, giọng không kém phần nhiệt tình: ``Tôi cũng thế\~{} Có gì ngon là tôi mang hết.''

Hai con người này đúng là\dots\ lúc nào cũng chỉ nghĩ đến ăn uống. Ryuuhou nhìn họ, rồi quay sang tôi, cười khúc khích: ``Em thấy hai chị này giống anh quá, lúc nào cũng nghĩ đến đồ ăn.''

Tôi lườm em: ``Anh không có thế. Cái này anh nói em mới đúng đấy.''

Ayame phì cười, còn tôi thì nhướng mày nhìn Korone và Okayu: ``Chỉ thế là nhanh thôi. Ít nhất thì hai đứa cũng phải phụ bọn này làm đi chứ? Không thể chỉ ăn không được.''

Chưa kịp để hai người kia phản ứng, một giọng nói khác vang lên từ phía sau, đầy vẻ tự hào: ``Chị biết trước là sẽ thế này nên đã chuẩn bị sẵn rồi.''

Fubuki vỗ hai nhịp, ngay lập tức một bộ dụng cụ mì trôi ống trúc xuất hiện như thể rơi xuống từ trên trời. Bộ dụng cụ này khá đơn giản: chỉ có một ống trôi xanh được khoét vừa để mì trôi xuống, vài chiếc giá đỡ ống cũng làm bằng ống trôi và một chiếc khay gỗ để đỡ mì nếu không gắp kịp. Trông nó như một món đồ thủ công được làm từ mấy ống tre đơn giản, nhưng cũng đủ cho một bữa tiệc mì trôi để đánh tan cái nóng oi ả.

Ryuuhou chạy lại, sờ vào ống trúc với vẻ mặt đầy thích thú: ``Fubuki-san giỏi thật! Làm sao chị có được cái này vậy?''

Nàng Cáo trắng cười, xoa đầu Ryuuhou: ``Chị có nguồn riêng mà. Bí mật nghề nghiệp.''

Chúng tôi không khỏi trầm trồ, còn Ayame thì cười khoái chí: ``Đúng là chúng em có thể tin tưởng vào chị mà! Có chị ở đây là mọi thứ sẽ ổn thôi.''

Ryuuhou cũng gật đầu lia lịa, như thể vừa tìm được một người chị lý tưởng: ``Em cũng muốn có một bộ như thế này! Có thể mua ở đâu không?''

Fubuki cười bí ẩn: ``Để chị mách cho sau.''

Cặp đôi Chó-Mèo không để lỡ thời gian. Mỗi người lấy ra một con tôm mà họ mang sẵn theo, rồi khẽ khàng động viên chúng như thể đang nói chuyện với thú cưng.

``Cố lên nhé, Tôm con!'' Okayu thủ thỉ, giọng như một người mẹ đang khích lệ con mình.

Korone cũng hào hứng không kém, đôi mắt sáng rực: ``Đi nào, Tôm gái! Chúng ta sẽ chiến thắng!''

\img{5-13c}

Tôi chớp mắt nhìn cảnh tượng trước mặt. Hai con tôm\dots\ vẫn còn sống và đang quẫy đạp. Chúng có màu đỏ ửng như đã bị luộc chín, nhưng vẫn giãy giụa mạnh mẽ. Thế quái nào họ tìm được hai con tôm đã áng đỏ mà vẫn còn sống được nhỉ?

Ryuuhou cũng há hốc mồm, chỉ tay vào hai con tôm: ``Ơ\dots\ hai chị lấy tôm ở đâu ra vậy? Trông chúng còn sống mà lại có màu đỏ-- tôm luộc chín mà còn bơi được à?''

Korone quay sang Ryuuhou, giọng như một chuyên gia động vật học: ``Đó là tôm biển sâu, em ạ. Chúng có màu đỏ tự nhiên, không cần luộc cũng vậy.''

Okayu gật đầu, nói thêm: ``Chị mua từ một cửa hàng đặc sản ở bến cảng á\~{} Chúng rất khỏe.''

Một dự cảm chẳng lành len lỏi trong đầu, nhưng trước khi kịp mở miệng hỏi xem hai người này thực sự lấy chúng từ đâu, cặp OkaKoro đã nhanh tay thả luôn hai sinh vật xấu số (tôi đoán vậy?) vào ống trúc. Hai con tôm rơi xuống dòng nước mát lạnh, và ngay lập tức, chúng lao đi với tốc độ tên lửa, để lại một vệt nước như đường băng siêu tốc. Tôi chỉ kịp thốt lên một câu đầy bất lực: ``Khoan đã! Hai người định làm gì với chúng thế!?''

Nhưng đã quá muộn. Kính cửa sổ phía xa nổ vỡ tan tành - lần thứ n rồi, và tôi cũng chẳng buồn phàn nàn nữa. Ryuuhou nhìn theo hai con tôm biến mất, rồi quay sang tôi với vẻ mặt đầy phấn khích: ``Anh ơi! Chúng chạy nhanh thật! Em muốn nuôi một con!''

Tôi đưa tay xoa trán, giọng đầy bất lực: ``Đừng có mà bắt chước mấy cái trò điên rồ đó. Tôm không phải để nuôi làm thú cưng.''

Hai con tôm tiếp tục phóng đi như thể đang tham gia một cuộc đua sinh tử, buộc chúng tôi phải đuổi theo đến tận bờ biển. Cả bọn lao ra khỏi văn phòng, chạy qua những con phố, và cuối cùng dừng lại trên bãi cát trắng mịn.

Korone vẫy tay đầy hào hứng, giọng như một cổ động viên cuồng nhiệt: ``Cố lên, Tôm gái! Tốc biến nào! Mày có thể làm được!''

Okayu cũng không chịu thua, giọng vẫn đều đều nhưng đầy khích lệ: ``Tôm con, đừng để bị bỏ lại phía sau đấy! Hãy cho thấy sức mạnh của mình!''

Fubuki với tốc độ siêu thanh của mình đã đáp xuống bãi biển (theo nghĩa đen) trước tiên. Khi chúng tôi đến nơi, em ấy đã cầm kính lúp xem xét tình trạng của hai con tôm vẫn còn sống nhăn răng sau màn đua không tưởng kia. Em ấy cúi xuống, nhìn chúng như một nhà khoa học đang quan sát một hiện tượng tự nhiên: ``Hmmm\dots\ Chúng vẫn còn sống. Thú vị thật.''

Chúng tôi dõi theo hai con tôm vẫn đang giãy đành đạch kia, và con tôm của Okayu (không có nơ) vượt trước đối thủ có đúng\dots\ vài xăng-ti-mét. Một khoảnh khắc căng thẳng, rồi Fubuki đứng dậy, giơ tay lên như một trọng tài: ``Chúc mừng, Tôm con của Okayu đã giành chiến thắng!'' Cô kéo tay Okayu lên cao, trong khi tôi và Ayame vẫn chưa thể hiểu chuyện quái gì đang diễn ra.

Ryuuhou nhìn hai con tôm, rồi nhìn Okayu, rồi lại nhìn tôi với ánh mắt đầy vẻ hâm mộ: ``Okayu-san giỏi quá! Tôm của chị chiến thắng rồi!''

Tôi lườm em: ``Có cái gì đâu mà reo hò ghê vậy\dots''

\begin{figure}[H]
  \centering
  \begin{subfigure}[t]{0.45\textwidth}
    \centering
    \includegraphics[width=8cm]{../../../../Image/Book5/5-13e1.png}
    \caption{Với kết quả suýt soát như vậy\dots}
  \end{subfigure}
  \quad
  \begin{subfigure}[t]{0.45\textwidth}
    \centering
    \includegraphics[width=8cm]{../../../../Image/Book5/5-13e2.png}
    \caption{\dots\ thì Okayu đã giành chiến thắng trong cuộc đua tôm này.}
  \end{subfigure}
\end{figure}

Korone rưng rưng nước mắt, cẩn thận nhấc con tôm của mình lên, giọng đầy tiếc nuối như một người mẹ vừa mất con: ``Tôm gái\dots\ Cưng cũng đã cố gắng hết sức rồi nhỉ. Chị tự hào về em.''

Sau đó, hai người họ cùng nhau thả hai con tôm về biển, vẫy tay chào tạm biệt như thể tiễn biệt một người bạn thân thiết. Ryuuhou cũng vẫy tay theo, giọng đầy cảm động: ``Tạm biệt hai chú tôm nhé! Hãy sống vui vẻ dưới biển nhé!''

``Đừng có lấy đồ ăn ra chơi chứ!'' Ayame thốt lên, bất bình với việc họ đem hai con tôm ra làm đồ chơi. ``Đó là tôm, không phải thú cưng!''

Korone lập tức quay lại với ánh mắt đầy hăm dọa, như một con chó đang bảo vệ lãnh thổ: ``Tôm gái là Tôm gái, không phải đồ ăn!''

\img{5-13f}

Tôi cười chát, giọng đầy bất lực: ``Nó còn sống nhăn răng thế kia thì đúng là không phải đồ ăn rồi\dots\ Cơ mà Ayame-chan nói đúng đó, đừng có lấy nó làm đồ chơi chứ. Lần sau muốn đua thì đua với cái gì khác đi.''

Ayame khoanh tay, lườm sang cặp đôi Chó-Mèo ấy: ``Thế thì đừng có ném nó ra ngoài chứ! Chúng có phải là đồ chơi đâu, yo!''

Ryuuhou chen vào, giọng đầy vẻ muốn hòa giải: ``Nhưng mà vui mà! Em chưa bao giờ thấy tôm đua bao giờ, phải công nhận là hay đấy.''

Tôi đưa mắt nhìn ra biển. Hai con tôm chắc giờ đã đoàn tụ với gia đình của chúng, tận hưởng những ngày tháng hạnh phúc trọn đời. Chắc thế. Hoặc cũng có thể tôi đọc quá nhiều truyện cổ tích rồi.

\img{5-13g}

Hoặc, nghĩ đến trường hợp tệ hơn, có khi chúng lại bị ai đó vớt lên rồi đem đi hấp tôm chín\dots\ ôi trời, tôi đang nghĩ cái quái gì vậy? Đây là \hll\ mà, cái quái gì cũng có thể xảy ra mà!

Ryuuhou đứng bên cạnh tôi, nhìn ra biển, giọng đầy mơ màng: ``Anh nghĩ tụi nó có sống hạnh phúc không?''

Tôi nhìn em, rồi thở dài: ``Anh hy vọng là có. Nhưng mà nếu chúng bị bắt lại, ít nhất thì chúng cũng đã có một cuộc đua đáng nhớ.''

\dots\ Thôi kệ. Quay trở lại vấn đề chính - bữa tiệc mì trôi ống trúc thực sự đang chờ chúng tôi, và tôi hy vọng rằng không còn ai mang thêm tôm ra nữa.

\ct

Sau khi chúng tôi thả hai con tôm về biển – với những lời chào tạm biệt đầy xúc động từ Korone và Okayu – Fubuki bỗng nảy ra một ý tưởng, đôi mắt cô sáng lên như vừa tìm thấy một kho báu: ``Đã ra đến đây rồi, sao ta không làm tiệc mì ống trúc ngay trên bãi biển luôn đi! Không khí biển, gió mát, và mì somen nguội – còn gì tuyệt hơn!''

Ayame và Okayu lập tức hưởng ứng với vẻ mặt hào hứng.

Ayame gật đầu, giọng đầy tán thành: ``Ý hay đấy. Biển với mì somen là combo hoàn hảo.''

Okayu cũng đồng tình, giọng vẫn đều đều nhưng có chút nhiệt tình: ``Triển luôn nào. Tôi đói rồi.''

Korone thì nhảy cẫng lên, hai tay giơ cao như một đứa trẻ vừa được hứa cho đồ chơi: ``Somen! Somen! Ăn thôi!''

Ryuuhou cũng hào hứng không kém, cô bé đứng cạnh tôi, nhìn ra biển với đôi mắt sáng rực: ``Em chưa bao giờ ăn mì somen trên biển! Nghe thú vị quá!''

Tôi chỉ chép miệng, nhún vai một cách cam chịu. Đằng nào chúng tôi cũng đã chạy đến đây rồi mà. Quay về văn phòng trong cái nóng này thì cũng chẳng khá hơn: ``Đành vậy thôi chứ biết sao được, chả lẽ quay về văn phòng giữa lúc đang vui? Mà\dots''

Tôi đột nhiên khựng lại, nhìn chằm chằm vào một thứ đang nổi lên từ làn nước xanh biếc. Một cái vây lớn, sắc nhọn, đang tiến thẳng về phía chúng tôi với tốc độ đáng sợ.

``\dots Mày ở đâu tòi lên vậy, cá mập kia?''

Một con cá mập trắng khổng lồ lao lên mặt nước, nhảy vọt khỏi mặt biển, miệng há rộng như đang chào đón chúng tôi. Tôi đứng hình, Ryuuhou thì hét lên một tiếng nhỏ, còn Ayame, với sự bình tĩnh đáng kinh ngạc, lập tức vung tay lên, nắm lấy chiếc chùy (hay cái gì đó tương tự) và hét lớn: ``\textbf{Cút ngay! Đây là bữa tiệc của bọn tao, không phải bữa trưa của mày!}''

Con cá mập, như thể hiểu được lời đe dọa, lặn xuống và biến mất ngay lập tức. Ryuuhou nhìn Ayame với ánh mắt đầy ngưỡng mộ: ``Ayame-sann cừ thật! Em chưa bao giờ thấy ai đuổi cá mập bằng cách nói chuyện cả.''

Ayame cười tự hào, vuốt nhẹ mái tóc dài: ``Kinh nghiệm sống đó, em ạ.''

Tôi chỉ biết lắc đầu cười trừ: ``Xem nạn nhân của vụ Sharknado năm ngoái nói gì kìa\dots''

Nhưng bầu không khí vui vẻ ngay lập tức bị phá tan khi Fubuki đột nhiên thốt lên thất thanh, người đờ ra như một bức tượng. Tôi và Korone đồng loạt quay sang hỏi, giọng đầy lo lắng: ``Sao thế? Có chuyện gì?''

Em ấy gục đầu xuống, giọng thất thần như một người vừa phát hiện ra mình quên mất thứ quan trọng nhất: ``Em quên mang mì theo rồi\dots\ Quên hết rồi\dots\ Làm sao mà tổ chức tiệc mì mà không có mì được chứ?''

Ryuuhou đưa tay lên trán, giọng đầy vẻ ``trời ơi'': ``Chị ấy quên mì thật à? Làm sao mà tổ chức tiệc mì mà không có mì được chứ?''

Sự im lặng bao trùm trong chốc lát trước khi cả bốn chúng tôi cùng buông một tiếng thở dài đầy thất vọng.

Ayame, Okayu và tôi đồng thanh ngán ngẩm, giọng như một bản hòa ca bất đắc dĩ: ``Ây chà chà\dots''

Khi cả nhóm đang chuẩn bị quay trở về với tâm trạng ngán ngẩm và hụt hẫng - tôi đang nghĩ đến việc gọi một chiếc xe taxi để về lại văn phòng - thì một giọng rao vang lên từ xa, vọng qua tiếng sóng biển và tiếng gió: ``\textbf{Có ai muốn ăn mì somen không\~{}?}''

Đó là Haato, trong bộ trang phục áo tắm sặc sỡ, đầu đội một chiếc nón rơm to tướng, tay đẩy một chiếc xe bán hàng nhỏ. Tôi không biết vì sao em ấy lại ở đây, trong bộ dạng này, nhưng tôi cũng chẳng thèm hỏi nữa.

\img{5-13i}

Nhóm chúng tôi lập tức quay đầu nhìn về phía giọng nói quen thuộc ấy và đồng thanh hô lên, giọng như một dàn hợp xướng bất ngờ: ``\textbf{Haachama-chama!}''

Ayame chạy tới, khuôn mặt hớn hở như vừa tìm được một tia hy vọng giữa bãi biển nắng cháy: ``Bọn em sẽ mua hết đống này luôn! Bao nhiêu cũng được!''

Haato cười rạng rỡ và nhanh chóng đáp lại, giọng đầy vẻ chuyên nghiệp: ``Có ngay! Một vắt là 200 yên, giá hợp lý cho một ngày nắng đẹp!''

Nhưng khi Haato nhìn thấy chúng tôi với bộ dụng cụ để ăn mì somen, cái ống trúc và khay gỗ mà Fubuki đã mang theo cùng với hộp lạnh đựng đá, đôi mắt cô sáng rực như thể vừa nhìn thấy một mỏ vàng.

\img{5-13j}

Ngay lập tức, em ấy thay đổi thái độ, hai mắt lấp lánh như thể đang nhìn thấy thần tài giáng trần. Giọng em ấy trở nên đầy tham lam, như một thương nhân đang nhìn thấy khách hàng giàu có: ``\textbf{Mười nghìn yên một vắt!}''

Ayame và tôi phản ứng lại ngay lập tức, giọng đầy bất bình:  ``Ăn cắp giữa ban ngày à!?'' / ``Sao đắt cắt cổ thế!? Đây là mì somen hay là vàng thế?''

Ryuuhou cũng nhảy vào, tay chỉ vào Haato với vẻ mặt đầy phẫn nộ: ``Chị định lừa tụi em hả? Mười nghìn yên một vắt thì em về nhà ăn cơm còn hơn!''

Haato vênh mặt lên, khoanh tay lại và hất hàm nói như một thương nhân lão luyện, giọng đầy tự tin: ``Cái đó là luật cung-cầu thôi, em à. Ở đây có mình tôi bán mì, các người đang đói, thế là đủ để tăng giá.''

Chúng tôi nhìn em ấy với sự ngán ngẩm, vội soát lại trong túi quần, nhưng rồi nhìn nhau với vẻ buồn thiu.

Ayame than thở, tay xoa bụng như một người đang đói đến phát điên: ``Tôi làm gì có nhiêu đó tiền\dots\ Tôi chỉ có mấy đồng xu lẻ trong túi thôi.''

Tôi cũng không khá khẩm hơn, lục soát túi quần và chỉ thấy một tờ giấy ghi chú cũ kỹ: ``Anh thì không mang tiền mặt với điện thoại đi. Ai lại mang tiền khi đi đuổi tôm chứ?''

Fubuki, Korone và Okayu cũng chịu chung số phận: họ đều mang tay không. Nhưng riêng cặp OkaKoro thì vẫn giữ nguyên vẻ mặt thản nhiên vui vẻ như chẳng có chuyện gì xảy ra. Họ đứng đó, mỉm cười, như thể đang chờ đợi một phép màu. Đám ăn chực đó thì hẳn nhiên làm gì có mang tiền trong người cơ chứ.

Ryuuhou nhìn cả nhóm, rồi quay sang tôi, giọng đầy thất vọng: ``Anh ơi, hay là mình về nhà ăn cơm trắng với muối vậy? Chứ em không có tiền đâu.''

Thấy cả nhóm chúng tôi không ai có tiền trong túi, Haato - với tư cách là một thành viên của \hll\ và cũng là một người có chút lương tâm - đành phải nhượng bộ. Em ấy hắng giọng, rồi miễn cưỡng hạ giá, giọng như đang đau đớn vì mất mát: ``V-Vậy thì năm nghìn yên cũng được! Đó là giá cuối cùng đấy!''

Nhưng cái giá mới mà cô nàng đưa ra vẫn chẳng khiến ai hài lòng. Cả năm người chúng tôi liếc nhìn nhau, trên mặt hiện rõ vẻ không chấp nhận được. Ryuuhou nhìn tôi, rồi nhìn Haato, rồi lại nhìn tôi, như thể đang truyền tải một thông điệp ngầm: ``Mặc cả tiếp đi.''

\img{5-13k}

Okayu là người lên tiếng trước, giọng bình thản như thể đang nói về thời tiết: ``Giá đó ai mua? Em có thể bán với giá nào đó hợp lý hơn không?''

Và thế là một cuộc mặc cả căng thẳng (nhưng hài hước) bắt đầu, với mỗi người một chiến thuật riêng.

Haato cố gắng giữ giá, giọng đầy quyết tâm: ``2000! Đó là giá rẻ nhất cho mì somen chất lượng cao!''

Korone lắc đầu, đôi tai cụp xuống như một chú chó bị từ chối: ``Không\~{} Giá đó thì em thà ra ngoài ăn hàng còn hơn.''

Haato nhăn mặt, giảm xuống một chút: ``Ư\~{} 1000! Đây là giá đặc biệt cho các bạn đấy!''

Fubuki đưa tay lên cằm, giọng đầy vẻ suy tư: ``Vẫn đắt quá đi\~{}. Em có thể bán với giá sinh viên được không?''

Haato thở dài, nhìn lên trời như thể đang cầu xin sự giúp đỡ: ``500! Cuối cùng! Đây là giá thấp nhất tôi có thể đưa ra!''

Ayame gật gù, nhưng vẫn chưa hài lòng: ``Gần được rồi! Nhưng vẫn còn cao.''

Haato gần như khóc, giọng tuyệt vọng như một thương nhân sắp phá sản: ``100 yên! Lấy 100 yên đi mà!''

Tôi nhún vai, với vẻ mặt của một người đã quen với việc mặc cả: ``Vẫn hơi đắt đấy. 50 yên thì may ra.''

Fubuki, Ayame, Korone và Okayu ngay lập tức gật đầu phụ họa theo tôi, tạo thành một dàn hợp xướng từ chối: ``Vẫn đắt quá!'' / ``Em có thể giảm thêm không?'' / ``Chúng tôi chỉ có mấy đồng xu thôi.''

Haato giật bắn người, giọng như sắp khóc: ``Em không thể giảm nữa đâu\dots\ Đây là giá cuối cùng rồi\dots\ Em sẽ lỗ vốn mất!''

Ryuuhou chen vào, giọng đầy vẻ thương lượng: ``Haato-san, nếu chị bán rẻ hơn, bọn em sẽ giới thiệu chị với bạn bè! Đó là quảng cáo miễn phí đấy!''

Tôi khoanh tay lại, nhún vai một cách bình thản, giọng như một người đã đạt được mục tiêu: ``Vậy thì\dots\ chúng tôi sẽ mua với giá 50 yên. Mỗi người một vắt, và coi như đó là một thỏa thuận.''

Cả nhóm chúng tôi đồng loạt cầm bát lên, bắt chước đúng giọng điệu ``cô bán hàng'' Haato khi nãy, giọng như một dàn hợp xướng vui vẻ: ``Ngon quá đi\~{}!''

Haato đờ người ra, hai tay xòe mấy đồng xu một trăm yên và vài đồng 50 yên như thể không tin vào thực tế phũ phàng này. Cô nhìn số tiền ít ỏi, rồi nhìn chúng tôi đang vui vẻ ăn mì, và thở dài: ``Lần sau em sẽ không bán cho mấy người nữa\dots''

Nhưng dù gì thì cuối cùng chúng tôi vẫn có một buổi tiệc mì trôi ống nước ngon miệng một cách miễn phí (gần như vậy). Ryuuhou vừa ăn vừa cười, giọng đầy hài lòng: ``Em thích món này! Lần sau lại tổ chức nữa nhé, Onii-san!''

Tôi nhìn em, rồi nhìn biển xanh và những người bạn đang cười đùa, và thầm nghĩ: ``\textit{Ít nhất thì cũng có một bữa tiệc đáng nhớ. Và lần sau, tôi sẽ nhớ mang tiền.}''

\begin{figure}[H]
  \centering
  \begin{subfigure}[t]{0.45\textwidth}
    \centering
    \includegraphics[width=8cm]{../../../../Image/Book5/5-13l1.png}
    \caption{Bốn cô gái (và tôi) có được một bữa tiệc mì somen\dots}
  \end{subfigure}
  \quad
  \begin{subfigure}[t]{0.45\textwidth}
    \centering
    \includegraphics[width=8cm]{../../../../Image/Book5/5-13l2.png}
    \caption{\dots trong khi Haato chỉ được nhận vài đồng bạc lẻ, coi như là động viên từ họ.}
  \end{subfigure}
\end{figure}

\pov{-}

\begin{center}
  \uline{Tất cả những gì diễn ra ở trên đều nằm trong đoạn phim.}
\end{center}

Đoạn phim kết thúc, và màn hình trở lại khung cảnh trong một căn phòng tràn ngập ánh sáng - văn phòng \hll\ quen thuộc. Ngồi trước màn hình là Kanata, đôi mắt vẫn còn dán chặt vào những gì vừa xem, vầng hào quang trên đầu khẽ lơ lửng như đang suy tư. Bên cạnh cô, Coco - người chịu trách nhiệm hướng dẫn buổi học hôm nay - gật đầu đắc ý, vỗ tay một cái rồi bắt đầu bài giảng của mình với giọng như một giáo viên đang truyền đạt kiến thức quý báu: ``Qua đoạn phim này, ta có thể rút ra một bài học quan trọng: hãy nhớ đừng bao giờ mua hàng từ bên trung gian nếu không muốn bị mua với giá cắt cổ! Như thế mới là một người tiêu dùng thông minh! Các bạn hãy ghi nhớ điều này khi đi chợ nhé.''

\img{5-13m}

Kanata gật gù theo phản xạ, nhưng trong đầu lại có một suy nghĩ hoàn toàn khác. Cô đưa tay lên cằm, nheo mắt nhìn màn hình đã tối màu: ``Cái tình huống này có vẻ hơi khiên cưỡng nhỉ\dots''

Tại sao tự nhiên bọn họ lại kéo nhau ra bãi biển, chỉ để rồi phát hiện ra là quên mất nguyên liệu chính? Rồi tự nhiên Haato xuất hiện với cả một xe mì somen như thể đã biết trước chuyện này? Và còn cả vụ thương lượng đầy kịch tính nữa chứ\dots\ Cô bắt đầu cảm thấy như mình đang xem một vở kịch được dàn dựng sẵn hơn là một bài học thực tế.

Kanata khoanh tay, nhướng mày nhìn Coco với ánh mắt đầy nghi ngờ của một học sinh đã bắt đầu hoài nghi giáo viên: ``Nhưng mà, sensei này\dots''

``Hửm?'' Coco quay sang, đôi mắt long lanh như một con Rồng đang chuẩn bị giải thích một bài học quan trọng.

``Tôi có cảm giác\dots\ cái đoạn phim này không phải là bài học về tiêu dùng, mà là về cách làm sao để\dots\ lấy đồ ăn miễn phí thì đúng hơn? Nhìn cách họ mặc cả kìa, rõ ràng là có chiến thuật hết.''

Coco nghe vậy suýt sặc, vội vã lấy lại bình tĩnh: ``N-Này! Đừng có suy nghĩ lệch lạc như thế chứ! Đó là một bài học về giá trị của đồng tiền và sự thông minh trong tiêu dùng!''

Kanata chỉ im lặng nhìn cô bằng ánh mắt đầy hoài nghi, như thể đang nhìn một kẻ lừa đảo vừa bị vạch mặt: ``Tôi lại thấy cái đoạn cả nhóm bắt chước điệu bộ của Haachama giống như một chiến thuật tâm lý hơn ấy. Họ đã dùng chính giọng điệu của cô ấy để hạ giá,đó là một đòn tâm lý chính hiệu.''

Coco đổ mồ hôi hột với lý luận của Kanata, cô đưa tay lau trán: ``Ờ thì\dots''

Cô ấp úng một chút rồi cố vớt vát, tay vung vẩy như thể đang cố gắng vẽ ra một bức tranh tươi sáng hơn: ``Nhưng ít nhất, bọn họ đã thành công trong việc tránh bị mua với giá cắt cổ, đúng không? Đó chính là bài học quan trọng: đôi khi, sự sáng tạo và tinh thần đồng đội có thể giúp bạn vượt qua những tình huống khó khăn về tài chính!''

Kanata thở dài, nhưng cũng không buồn tranh luận thêm. Cô biết rằng với Coco, mọi thứ đều có thể trở thành một bài học về kinh tế học, ngay cả một bữa tiệc mì somen trên bãi biển: ``Vâng vâng, tôi hiểu rồi\dots\ nhưng mà\dots\ chúng ta có thể học những bài học thực tế hơn không? Ví dụ như cách nấu mì somen ngon, hay cách tránh bị cá mập tấn công khi đang ăn trên biển?''

Coco cười trừ, ho khẽ một tiếng rồi vội vàng chuyển sang bài giảng tiếp theo, như thể đang cố gắng né tránh một cuộc chất vấn khó khăn: ``Được rồi, vậy tiếp theo chúng ta sẽ học về cách đàm phán khi đi chợ\dots\ Với các tình huống thực tế, và không có cá mập.''

Kanata bất giác cảm thấy hơi lo lắng. Không lẽ bài học tiếp theo lại là một đoạn phim có sự góp mặt của Haachama nữa sao\dots? Cô nhìn Coco với ánh mắt đầy nghi ngờ, nhưng không nói gì thêm, chỉ thở dài và chuẩn bị cho bài học tiếp theo.

%==========================================TRUYỆN PHỤ=========================================
\omk

\pov{Hikawa Kuon}

\begin{center}
  \uline{Hậu trường ghi hình}
\end{center}

``Và\dots\ xong!''

Chúng tôi đã ghi hình xong thước phim giáo dục. Không giống lần trước khi Miko là người phụ trách - và cô ấy đã làm một video về an toàn lái xe mà cuối cùng trở thành một vở hài kịch - lần này Coco sẽ đảm nhận vai trò tuyên truyền về ``mua hàng qua trung gian''.

\dots Tại sao lại họ muốn làm theo một chủ đề này chứ? Cái bộ phim giáo dục về an toàn lái xe năm ngoái\footnote{Xem lại {\flaregothic ÉPISODE 14}, quyển số Một.} họ còn làm chẳng ra gì mà! Tôi đã thấy Coco ghi hình và cô ấy đã cười suốt cả buổi, đến nỗi tôi không biết đây là một video tuyên truyền hay một chương trình tạp kỹ rẻ tiền nữa.

``Tại sao mấy người lại cho tôi đóng vai gian thương chứ!?'' Haato khoanh tay, mặt mày bực bội, đôi má phồng lên như một đứa trẻ bị bắt nạt. ``Em là một idol trong sáng, chứ không phải một kẻ bán hàng rong với giá cắt cổ!''

Cả nhóm chỉ nhìn nhau cười khúc khích. Ryuuhou đứng bên cạnh tôi, cười đến mức phải vịn vào vai tôi để khỏi ngã: ``Haato-san đóng vai gian thương hợp quá trời hợp luôn! Em nghĩ chị ấy nên theo đuổi sự nghiệp diễn xuất đấy.''

Fubuki bụm miệng cười, giọng như một người đang cố gắng giữ bí mật nhưng không thể: ``Thì mặt bà gian xảo mà. Trông bà lúc mặc cả y hệt một thương nhân chuyên nghiệp ngoài chợ.''

Ngay sau đó, tôi cũng cất giọng châm chọc, nhớ lại những lần Haato đã bày trò trêu chọc cả nhóm: ``Và cũng là người bày trò trêu mọi người nhiều nhất đấy. Em đã từng giấu chuột máy tính của Coco trong tủ lạnh, và bây giờ em đóng vai một kẻ bán hàng lừa đảo, quá hợp lý còn gì.''

Haato nổi sần sần lại, mặt đỏ bừng như bị chúng tôi gãi đúng chỗ ngứa: ``\textbf{Là sao hả!? Em chỉ đùa thôi mà! Đó là những trò đùa vô hại!}''

Ryuuhou cười lớn, vỗ tay: ``Chị mà gọi đó là vô hại á? Em nhớ lần chị đổ nước vào giày của Korone-san đấy.''

Tất cả bật cười trước phản ứng của nàng Trung học. Cô nàng phồng má, hờn dỗi, nhưng không thể chối cãi rằng hình tượng của mình thực sự rất hợp vai một kẻ bán hàng gian xảo với giá cắt cổ.

Nhưng tôi thì vẫn thấy có một điều gì đó chưa ổn. Tôi đưa tay lên cằm, nhìn lại đoạn phim trên màn hình máy quay: ``Nhưng mà anh thấy phần này vẫn chưa ổn\dots\ Cái đoạn mà chúng mình mặc cả ấy. Nó chưa có sự kịch tính cần thiết.''

Haato cũng lên tiếng bất bình với vai của em ấy, tay chỉ về phía màn hình: ``Đúng đó! Bình thường giá đắt là người ta sẽ nói `Không mua thì cút!' kiểu thế\dots\ Đây lại mềm lòng giảm giá, thế có chết không!? Em đã phải đóng vai một kẻ bán hàng mà không được chửi ai, thật bất công!''

Tôi khoanh tay, nhớ về những ngày mà mẹ tôi hay trả giá mỗi khi mua đồ ở ngoài chợ: ``Anh không biết bọn em thế nào, nhưng ở chỗ bọn anh, mấy đứa chỉ cần thêm cái quay ngoắt đi lạnh lùng, không thèm đoái hoài nữa là nó sẽ hợp lý hóa phân cảnh vừa rồi ngay. Thế mới là mặc cả chuyên nghiệp.''

Ryuuhou nhìn tôi với vẻ ngạc nhiên, như thể vừa phát hiện ra một bí mật: ``Anh biết nhiều về chuyện này thế? Hay là anh đã từng đi chợ cùng mẹ nhiều?''

Tôi nhún vai, giọng đầy vẻ tự hào: ``Kinh nghiệm sống đó em ạ.''

Korone nhìn tôi với sự ngạc nhiên, đôi tai chó dựng lên: ``Thế á? Em chưa bao giờ nghĩ đến việc đó.''

Okayu cũng nói góp, giọng vẫn đều đều nhưng có chút tò mò: ``Bọn em không biết đấy. Ở nhà bọn em toàn ăn đồ có sẵn, không cần mặc cả.''

Cuộc bàn luận về đoạn phim vẫn chưa đến hồi kết thì Ayame lên tiếng, cắt ngang câu chuyện của chúng tôi với vẻ mặt đói khổ của một người vừa trải qua một ngày dài: ``Ờm\dots\ Mọi người có muốn ăn tiếp không? Em thấy vẫn còn mì somen trong xe của Haato-senpai kìa.''

Không cần suy nghĩ, tất cả đều đồng thanh hô lớn, giọng như một dàn hợp xướng vui vẻ: ``Tất nhiên là có chứ!''

Ryuuhou nhảy cẫng lên, giơ hai tay lên trời: ``Em đói rồi! Ăn tiếp đi!''

Tôi quay sang hỏi nàng Song tính, giọng đầy thiện chí: ``Haato-chan có muốn ăn với bọn này không? Coi như là em đã mời chúng tôi, và chúng tôi mời em.''

Haato thoáng do dự, nhưng rồi bày ra vẻ mặt e dè, như thể đang sợ bị từ chối: ``Nhưng mà còn bát với đũa không? Em không muốn ăn bằng tay đâu.''

Tôi nhún vai, cười nhẹ, chỉ về phía chiếc túi đựng dụng cụ của Fubuki: ``Chuẩn bị sẵn một bộ, đề phòng ai đó đánh rơi. Anh đã nghĩ đến điều này từ đầu rồi.''

Haato bật cười, rồi hào hứng ngồi xuống cùng chúng tôi, giọng như một đứa trẻ vừa được nhận quà: ``Vậy thì em không khệ nệ mọi người nha\~{}!''

Một trận mì somen trôi ống nước nữa lại diễn ra, và ai cũng có phần của riêng mình. Ryuuhou ngồi bên cạnh tôi, vừa gắp mì vừa nói: ``Onii-san, mì somen này ngon thật! Lần sau em muốn học cách làm, để về nhà làm cho bố mẹ ăn.''

Tôi nhìn em, giọng đầy vẻ bất ngờ: ``Em biết nấu ăn à?''

Ryuuhou cười tinh quái: ``Em học được từ Y*uT*be đấy. Có gì khó đâu.''

Tôi là người ăn ít nhất, vì một thằng đàn ông như tôi sẽ phải nghĩa vụ nhường phần cho các cô gái. Cô em gái thấy vậy. nhìn tôi với vẻ không hài lòng: ``Anh ăn ít thế? Em thấy anh cứ gắp mấy sợi mì rồi thôi à.''

Tôi đáp, giọng đầy vẻ khiêm tốn: ``Anh không đói lắm. Mà các em ăn đi, kẻo nguội.''

\dots\ Mặc dù sau đó họ cũng đã cố thuyết phục tôi rằng: ``Cứ ăn nhiều vào, Kuon-senpai! Năm người bọn em ăn không nổi cái xe mì này đâu!''

Cuối cùng tôi cũng phải ăn thêm một bát nữa vì bị cô em gái lém lỉnh và nàng Oni lắm mưu ấy ép.

Sau khi ăn xong, mọi người rủ nhau xuống biển chơi. Họ cũng rủ tôi đi cùng, nhưng tôi từ chối, lấy lý do cần có người trông đồ. Thành thật mà nói, tôi không thích tắm biển lắm, nhất là khi không mang theo đồ bơi. Ryuuhou nhìn tôi, giọng đầy vẻ tiếc nuối: ``Anh không đi à? Biển đẹp lắm đấy.''

Tôi xoa đầu em, giọng ấm áp: ``Anh ở đây trông đồ. Các em đi chơi đi, đừng lo.''

Cả nhóm vui chơi đến tận chiều tối. Tiếng cười, tiếng sóng biển, và những bước chân chạy trên cát tạo nên một bản hòa ca mùa hè hoàn hảo. Khi quay về, ai cũng trông rạng rỡ, tràn đầy năng lượng sau một ngày quậy phá. Ryuuhou chạy lại, tóc vẫn còn ướt, mặt đỏ bừng vì nắng: `Onii-san! Biển ở đây mát quá! Lần sau em lại đi nữa nhé!''

Tôi nhìn em, không thể không mỉm cười: ``Được rồi, nhưng nhớ mang đồ bơi nhé.''

Trước khi tạm biệt nhau, tôi hứa với họ sẽ tìm cách xây một bộ mì trôi ống trúc mà tôi từng thấy trên mạng: một bộ to hơn, chắc chắn hơn, và có thể dùng được nhiều lần. Ngay lập tức, tất cả đều háo hức mong chờ, nhất là Okayu và Korone - hai con người thích ăn ``chùa.'' Ayame, Fubuki và Haato cũng cười rạng rỡ, hăng hái đồng ý.

Ryuuhou nắm tay tôi, giọng đầy hứa hẹn: ``Em sẽ giúp anh làm! Em biết cách cưa tre đấy!''

Tôi nhìn em, giọng vừa bất ngờ vừa lo lắng: ``Em biết cưa tre à? Đừng có mà đứt tay đấy.''

Một mùa hè nữa lại khép lại, cùng với đó là những kỷ niệm mới mà tôi và \hll\ đã dành cho nhau. Những tiếng cười, những cuộc mặc cả, và những bữa tiệc mì somen trên bãi biển, tất cả đều trở thành những mảnh ghép đẹp đẽ trong bức tranh mùa hè của chúng tôi.

Cơ mà\dots\ vẫn còn một ngày nữa mới hết tháng Tám nhỉ?

Tôi có cảm giác như\dots\ chúng ta vẫn còn một điều gì đó chưa xảy ra\dots\ Một linh cảm mơ hồ, như thể có một điều gì đó đang chờ đợi chúng tôi ở phía trước.

Tôi nhìn về phía hoàng hôn, nơi mặt trời đang dần khuất sau những đám mây, và thầm nghĩ: ``\textit{Mình không biết đó là gì, nhưng chắc chắn nó có dính dáng đến đồ ăn. Có thể là một bữa tiệc khác\dots\ hoặc có thể là một món cà ri nào đó.}''

Ryuuhou kéo tay tôi, giọng đầy hào hứng: ``Anh đang nghĩ gì thế, Onii-san?''

Tôi nhìn em, rồi nhìn những người bạn đang vẫy tay từ xa, và mỉm cười: ``Cuối tuần này Choco-san đang định dựng một thước phim hướng dẫn nấu món cà-ri. Em có muốn tham gia không?''

Ryuuhou nhảy cẫng lên, giọng như một đứa trẻ vừa được hứa tặng quà: ``Có chứ! Em sẽ mang cả bố mẹ đến nếu được!''

Tôi phì cười, xoa đầu em: ``Đừng có mang cả nhà đến đấy. Nhưng mà\dots\ có lẽ một bữa tiệc cà ri sẽ là một kết thúc hoàn hảo cho mùa hè này.''

\end{document}